\bookchapter{第十九章：故障排查與自救}

\chapterintro{手機如人，偶爾亦會“生病”。卡頓、閃退、黑屏、充不進電……大部分毛病實則未至那般可怕，君完全可於送修之前己先試之。此章乃君之“手機急診手冊”——先以低成本之法試，不行再覓專業醫生。}

\booksection{一、手機變卡、發熱——最常被問之問題}

實言之，手機用一年左右始變卡，太正常矣。原因未必乃手機壞矣，更多時候乃“積灰”——垃圾文件、後臺進程、存儲空間慢慢被填滿。

\booksubsection{先作診斷，勿急修：}

\begin{enumerate}
  \item \textbf{查存儲空間}：往設置 存儲視剩餘空間。若不到10\%矣，系統將變得甚卡，因其掙扎着騰地方。吾之經驗乃：保持至少15\%-20\%之剩餘空間，手機方能較順暢。
  \item \textbf{查耗電排行}：往電池設置中視過去24時之耗電排行。若某個App佔了大頭然君根本未甚用，即其在後臺偷跑。直卸載或限後臺活動。
  \item \textbf{聯想近來做了何}：乃升級系統後變卡？（大版本更新後前幾日變卡正常，系統於後臺重建索引）。抑或裝了某個新App之後變卡？（卸載試視）。
  \item \textbf{視使用環境}：手機於太陽底下暴曬、套着厚厚之殼邊充邊玩遊戲，發燙乃必然者。高溫尚將致CPU自動降頻，君即覺手機突然變卡矣。涼下來再試。
\end{enumerate}

\textbf{治療之法（自簡至繁）}：重啓手機（解六成之問題） 清存儲空間（刪無用之App與視頻） 某個App搞鬼即清其緩存或卸載重裝 最後一招：備份數據後恢復出廠設置。

\booksection{二、應用閃退、無響應——App不聽話矣}

某個App點開即閃退、或用着用着突然卡死不動矣。處理流程如醫生問診，一步步來：

\begin{enumerate}
  \item \textbf{先關重開}：自後臺劃掉App，再重新開。大多數臨時bug此下即好矣。
  \item \textbf{更新App}：往應用商店視有無更新。有時乃因App版本太老，與系統不兼容。
  \item \textbf{查權限}：有些App需某個權限方能正常工作，然君未給。比如相機App未給相機權限開了即閃退。往設置中覓彼App，將該給之權限皆給了。
  \item \textbf{清緩存}：往設置 應用管理 覓彼App 存儲 清除緩存。注意乃"清除緩存"非"清除數據"！清除數據將把君App內之登錄狀態與設置全刪了，相當於恢復出廠給App之狀態。
  \item \textbf{卸載重裝}：最後一招。卸載之前先確認此App之數據會否丟——微信之話先備份聊天記錄，遊戲之話先確認賬號有否綁手機號。
\end{enumerate}

\booksection{三、充不進電或充電慢——恐非電池之鍋}

充不進電乃尤易令人焦慮之問題，然甚多時候原因出於配件上，非手機壞矣。

\begin{itemize}
  \item \textbf{先排查充電口}：以牙籤或不導電之小工具，輕輕將充電口內之灰塵與棉絮掏之。此吾己遇過——口袋裏掉了一根線頭剛好堵於充電口內，插進去以爲接觸不良，掏出即好矣。
  \item \textbf{換根線與充電頭試}：線之損壞率比充電頭高甚多，尤其經常捲來捲去之線。換一條確定沒問題之線與充電頭交叉測之，很快能定乃哪個配件壞矣。
  \item \textbf{無線充電之問題}：若君用無線充電，略查手機殼是否太厚，或殼後面有無金屬片（磁吸支架之類）。又，將手機準確地放至充電板之正中——歪了即或充不上。
  \item \textbf{充電慢}：確認充電頭功率夠不夠（iPhone欲快充至少需20W之PD充電頭），線支不支持快充（非所有線皆能過3A以上電流），手機是否太燙（太燙將自動降功率護電池）。
\end{itemize}

\booksection{四、網絡連不上、Wi-Fi老是斷}

於外面用流量好好的，歸家連不上Wi-Fi；或Wi-Fi連着然就是沒網。按此序排查：

\begin{enumerate}
  \item \textbf{開關飛行模式}：下拉控制中心，開飛行模式等5秒再關之。此等於令手機重新搜一次網，甚有效。
  \item \textbf{Wi-Fi"忘記"再重連}：往Wi-Fi設置中，點彼網絡名旁之感嘆號（安卓乃齒輪），選"略之此網絡"或"忘記"，而後重新輸密碼連之。
  \item \textbf{重啓路由器}：多人家裏之路由器數月未關過，時久將不太穩定。拔掉電源等30秒再插回去，大部分Wi-Fi問題即解矣。
  \item \textbf{查是否手機之問題}：別的設備連同一Wi-Fi正常否？若都正常即君之手機不行，則手機或需重啓或重置網絡設置（設置 通用 傳輸或還原 還原網絡設置，不會刪君之數據，然將把所有Wi-Fi密碼皆清矣）。
\end{enumerate}

\booksection{五、屏幕問題：觸控不靈、閃屏、留印子}

\begin{itemize}
  \item \textbf{觸控不靈敏}：先撕掉鋼化膜視是否膜之鍋——便宜鋼化膜用久確實將影響觸控。屏上有水或手上有油亦將致觸控漂移。皆非之話，略重啓再試。
  \item \textbf{屏上有印子/殘影（燒屏）}：長時間顯同一畫面（比如一直亮屏刷抖音），將於OLED屏上留"殘影"。此非屏壞矣，只暫時的。關屏休息半個時一般能恢復。
  \item \textbf{屏閃爍}：若是低亮度下閃爍，恐乃PWM調光之問題，多見於OLED屏。往設置中覓"DC調光"或"防頻閃"之項，開後將佳多。非所有手機皆有此能。
  \item \textbf{屏有一塊區域全黑/亮}：此較嚴重，大概率乃屏物理損壞，需換屏。非官方維修換之屏，色彩與亮度通常不如原裝，亦影響防水。
\end{itemize}

\booksection{六、無法開機、卡於Logo界面——先勿慌}

此情最易令人心跳加速，然冷靜，先試此數個：

\begin{itemize}
  \item \textbf{充電半時再說}：甚多時候只是電徹底耗盡矣。插上原裝充電器，等至少30分鐘，再嘗試開機。有些手機於電耗盡後需充至一定電量方會有反應。
  \item \textbf{強制重啓}：iPhone全面屏機型按一次音量加、一次音量減、長按側邊電源鍵不放，直至屏現蘋果Logo（或要按住10-15秒）。安卓通用之法：長按電源鍵10秒以上，直至震動或亮屏。
  \item \textbf{卡於啓動界面反覆重啓}：此情叫"白蘋果"（iPhone）或"卡Logo"（安卓）。若強制重啓不行，或需入恢復模式以電腦修復。iPhone連電腦開iTunes/Finder，入恢復模式後選"更"（非"恢復"），將重裝系統然保留數據。安卓進Fastboot或Recovery模式，各品牌組合鍵不同，議搜君手機型號+"進恢復模式"。
  \item \textbf{以上皆不行}：即只能求助專業維修矣。數據若備份過，即不用太憂。
\end{itemize}

\booksection{七、聲音問題：喇叭不響、通話聲音小}

\begin{itemize}
  \item \textbf{喇叭無聲}：先確定非靜音模式（iPhone側邊之靜音撥片）、勿擾模式、或音量被按至最低矣。聽來甚傻，然此情佔了四成。而後查藍牙是否連着某個耳機或音箱。
  \item \textbf{通話對方聽不到君言}：查話筒（麥克風）孔有無堵住——手機底部之小孔甚易被灰塵堵住。以軟毛刷輕輕清理。勿以針或牙籤硬捅，將捅壞防塵網。
  \item \textbf{聲沙啞/破音}：若是喇叭進水所致——勿以吹風機熱風吹！以紙巾吸乾表面水分，而後放着自然晾乾。網上言用聲波排水App有一定之效（播放低頻聲把水振出），可試之。
\end{itemize}

\booksection{八、覓官方維修，避彼等坑}

若以上之法皆試過猶不行，即須覓專業維修矣。然維修市場水甚深，述吾己踩過之坑與經驗：

\begin{itemize}
  \item \textbf{覓真官方，勿覓"看起來像官方"者}：往品牌官網（Apple.com.cn、Huawei.com.cn等）覓"授權服務商"列表。街邊掛個蘋果Logo之非授權店，用之或非原裝配件，修完君之防水即無矣。
  \item \textbf{送修前必做之事}：備份所有數據（勿指望修之時數據不會丟），退所有要緊賬號（Apple ID、微信、支付寶），關"查找吾之手機"。有些維修點將要求君提供鎖屏密碼，此事有隱私之險——君最好己先退賬號再送過去。
  \item \textbf{換屏之坑}：非官方屏分"原壓屏"（原裝玻璃換過）、"後壓屏"（非原裝玻璃）、"組裝屏"（全非原裝）。價差好幾倍，然顯示效與耐久度亦差甚多。若手機尚在保修期，老老實實走官方。
  \item \textbf{換電池之坑}：非原裝電池或虛標容量，且用半年即衰減得比原裝電池二年之損還大。蘋果換電池官方價比非官方貴不了太多，議走官方。
  \item \textbf{數據恢復之坑}：手機進水/摔壞後，往維修若對方言"數據恢復要加錢"，小心——有些不良商家將故意誇大損壞程度以多收錢。正規之數據恢復乃先報價再操作，且恢復不成功不收費。
\end{itemize}

\begin{tcbwarn}{萬能第一招：重啓手機}
萬能第一招：重啓手機

\end{tcbwarn}

\begin{tcbpractice}{今即記之：此等救命之操作當爛熟於心}
\begin{enumerate}
  \item \textbf{強制重啓}：全面屏iPhone按音量加 音量減 長按電源鍵至現Logo。安卓長按電源鍵10-15秒直至震動。
  \item \textbf{清理充電口}：覓一小工具（牙籤或手機取卡針），輕輕清理充電口之灰塵。
  \item \textbf{Wi-Fi重連}：往Wi-Fi設置中"忘記"君當前之Wi-Fi，重新輸密碼連之。
  \item \textbf{開/關飛行模式}：下拉控制中心，開飛行模式，心中數5秒，關之，視網絡有否恢復。
  \item ：下拉控制中心，開飛行模式，心中數5秒，關之，視網絡有否恢復。
\end{enumerate}

\end{tcbpractice}
